\chapter{Conclusion} \label{Chap6}

In scenarios with sparse rewards or unreliable reward functions, the choice of the RL algorithm directly affects the safety and reliability of the policy. Through systematic experiments on the HumanoidBench Door task, this thesis reveals fundamental differences between PPO and GRPO when facing such reward conditions:

\begin{enumerate}
  \item \textbf{PPO engages in reward hacking when rewards are unreliable.} When the dense reward is not fully aligned with the true objective, PPO's return soars from a baseline of 244 to 431 ($+77\%$), while the success rate drops from 5.0\% to 0.0\%; its GAE advantage estimate amplifies the reward components weakly related to the success condition and suppresses the component related to the true goal. After removing the unreliable dense reward, PPO's policy completely collapses, with the return falling below 52\% of the baseline level.

  \item \textbf{GRPO remains stable under both conditions, and its stability stems from KL-return decoupling rather than policy stagnation.} GRPO's return stays within a narrow band of 188--219 across all seeds; the update magnitude (KL 3.8--25.2 nats, parameter distance 0.8--2.4) varies with seed and condition, but the return does not drift with the update magnitude. Under extremely sparse rewards, GRPO maintains the same exploration level as the Baseline ($\log\sigma$ shift of only $+0.002$) and avoids deterministic collapse. In contrast, PPO's KL and return are tightly coupled: under the Misaligned condition, higher KL means more severe reward hacking; under the Sparse condition, it exhibits two failure modes, namely KL-explosion deterministic collapse and KL-collapse gradient vanishing.

  \item \textbf{Reward decomposition and KL-penalty ablation jointly diagnose the source of stability.} Component-wise reward decomposition shows that PPO's return gain is concentrated in components weakly related to success such as door openness, while the passing-the-door component that truly reflects task completion decreases; the KL-penalty ablation rules out the alternative explanation that ``stability stems from the absence of the KL penalty'' (Section~4.3). Together these lines of evidence point to the normalization effect of group-relative advantage: components that are ``consistent across the group'' do not generate an advantage signal, so GRPO does not systematically amplify any single reward component.

  \item \textbf{Cross-task evidence and practical implications.} Additional experiments on the Reach task show that the conclusions are not limited to the Door task. In scenarios where reward design is imperfect or rewards are extremely sparse, we recommend GRPO as the fine-tuning starting point: its robustness to the reward structure does not depend on the setting of the KL penalty, and its performance lower bound is significantly better than PPO's (Section~4.1). Model selection should jointly monitor reward component decomposition and success rate rather than return increases alone; in our experiments the 77\% return growth of PPO was precisely a misleading signal for selection, while selecting checkpoints by success rate yields conclusions consistent with the main experiments (Section~4.1). The applicability boundary of this recommendation is: when the reward completely lacks intermediate discrimination, such as the purely binary reward of the Reach task, the within-group normalization of group-relative advantage degenerates into noise and GRPO's advantage narrows (Section~4.4).
\end{enumerate}
